\documentclass[12pt]{article}

% Idioma e codificação
\usepackage[brazil]{babel}
\usepackage[utf8]{inputenc}
\usepackage[T1]{fontenc}

% Matemática e listas
\usepackage{amsmath, amssymb}
\usepackage{enumitem}

% Layout e cores
\usepackage{geometry}
\usepackage{graphicx}
\usepackage{xcolor}

% Cores institucionais
\definecolor{maincolor}{RGB}{0,100,50}
\definecolor{forestgreen}{RGB}{0,100,50}

% Cabeçalho e rodapé
\usepackage{fancyhdr}

% Blocos
\usepackage[most]{tcolorbox}
\newtcolorbox{block}[1]{%
  colback=green!5,
  colframe=maincolor,
  coltitle=white,
  title=#1,
  fonttitle=\bfseries,
  boxrule=1.5pt,
  arc=3mm,
  left=6pt,
  right=6pt,
  top=6pt,
  bottom=6pt
}

% Margens
\geometry{margin=2.5cm}

% Fancyhdr
\pagestyle{fancy}
\fancyhf{}

\renewcommand{\headrulewidth}{2pt}
\renewcommand{\footrulewidth}{2pt}

\renewcommand{\headrule}{\hbox to\headwidth{%
  \color{forestgreen}\leaders\hrule height \headrulewidth\hfill}}

\renewcommand{\footrule}{\hbox to\headwidth{%
  \color{forestgreen}\leaders\hrule height \footrulewidth\hfill}}

% Altura do cabeçalho
\setlength{\headheight}{52pt}
\addtolength{\topmargin}{-20pt}

% Cabeçalho
\fancyhead[L]{%
  % Substitua pelo caminho correto da imagem ou mantenha o texto se não houver arquivo
  \textbf{UniRV} \\ \small Universidade de Rio Verde
}

\fancyhead[R]{%
  \raggedleft
  \textcolor{forestgreen}{%
    \small
    \textbf{Lista de Exercícios: Métricas Google SRE}\\
    Me.~Sergio Souza Novak\\
    \textit{Engenharia da Qualidade e Confiabilidade}
    }
}
    
% Rodapé
\fancyfoot[C]{\textcolor{forestgreen}{\thepage}}

% Comando para caixas de resposta
\newtcolorbox{resposta}[1]{%
  colback=white,
  colframe=forestgreen!30,
  arc=1mm,
  height=#1,
  width=\textwidth,
  boxrule=0.8pt,
  breakable
}

% =========================
\begin{document}
% =========================

\begin{center}
  {\Large \textbf{Lista de Exercícios: Métricas Google SRE}}\\[0.2cm]
  \textit{Engenharia de Software / Engenharia da Qualidade}
\end{center}

\vspace{0.4cm}

\section*{Parte 6 — Métricas de Requisição e Confiabilidade em APIs}

\textbf{16.} Explique a diferença entre disponibilidade baseada em tempo e disponibilidade baseada em requisições.
\begin{resposta}{3.5cm}
\end{resposta}

\textbf{17.} Em uma API REST, por que erros 4xx normalmente não são contabilizados como falhas do sistema para cálculo de disponibilidade?
\begin{resposta}{3.5cm}
\end{resposta}

\textbf{18.} Um sistema recebeu 1.200.000 requisições em um dia. 
Foram registrados:
\begin{itemize}
    \item 1.199.400 respostas 2xx
    \item 400 respostas 4xx
    \item 200 respostas 5xx
\end{itemize}

\begin{enumerate}[label=\alph*)]
    \item Calcule a disponibilidade considerando apenas 2xx como sucesso.
    \item Calcule a disponibilidade considerando apenas 5xx como falha real do sistema.
\end{enumerate}

\begin{resposta}{5cm}
\end{resposta}

\textbf{19.} Um serviço possui meta de disponibilidade de $99,9\%$ ao dia. 
Em um total de 2.500.000 requisições diárias:
\begin{enumerate}[label=\alph*)]
    \item Quantas falhas são permitidas por dia?
    \item Se ocorrerem 1.200 erros 5xx, a meta foi cumprida?
\end{enumerate}

\begin{resposta}{5cm}
\end{resposta}

\textbf{20.} Durante um teste manual de QA foram realizadas 6 requisições:
\begin{itemize}
    \item 3 retornaram 2xx
    \item 2 retornaram 4xx
    \item 1 retornou 5xx
\end{itemize}

\begin{enumerate}[label=\alph*)]
    \item Calcule a disponibilidade considerando apenas 2xx.
    \item Calcule considerando apenas 5xx como falha do sistema.
    \item Explique qual abordagem é mais adequada em ambientes de produção.
\end{enumerate}

\begin{resposta}{6cm}
\end{resposta}

\textbf{21.} Defina POFOD aplicado a uma API HTTP. 
Se um serviço apresenta 350 falhas 5xx em 700.000 requisições, calcule o POFOD.
\begin{resposta}{3.5cm}
\end{resposta}

\textbf{22.} Um sistema apresenta ROCOF = 0,002 falhas por hora.
\begin{enumerate}[label=\alph*)]
    \item Calcule o MTTF.
    \item Se o MTTR médio é 30 minutos, calcule a disponibilidade.
\end{enumerate}

\begin{resposta}{5cm}
\end{resposta}

\textbf{23.} Explique o conceito de Error Budget e sua relação com metas de 99,9\%, 99,99\% e 99,999\%.
\begin{resposta}{4cm}
\end{resposta}

\textbf{24.} Em um sistema de monitoramento industrial, o que é mais crítico:
Disponibilidade ou Confiabilidade? Justifique considerando riscos operacionais.
\begin{resposta}{4cm}
\end{resposta}

% =========================
\end{document}
% =========================